% 第9章：黑洞理论 - 包含阶段A/B/C/D完整内容
\section{黑洞理论}
\label{sec:bh}

\subsection{引言：黑洞的三大谜题}

黑洞对量子引力提出最严峻的挑战：
\begin{itemize}
    \item \textbf{信息悖论}：落入黑洞的量子信息会发生什么？霍金辐射表现为热的，暗示信息丢失
    \item \textbf{奇点问题}：广义相对论在 $r = 0$ 处预言无限密度，物理定律崩溃
    \item \textbf{熵问题}：Bekenstein-Hawking熵 $S_{BH} = A/4G$ 暗示几何、热力学与量子信息之间的深刻联系
\end{itemize}

CTUFT通过分形-扭转几何和10维内空间解决所有三个问题。

%==============================================================================
\subsection{阶段A：扭转场的正则量子化}
%==============================================================================
\label{subsec:stageA_canonical}

\subsubsection{无限维辛几何框架}

对于场论系统，相空间是\textbf{无限维辛流形}。设扭转场为 $\mathcal{T}^\alpha_{\mu\nu}(x)$，其中 $\alpha$ 是内部指标，$x \in \Sigma$（空间切片）。

\textbf{辛形式的构造}：

扭转场的相空间由场变量 $\{\mathcal{T}^\alpha_{\mu\nu}(x), \pi_\alpha^{\mu\nu}(x)\}$ 构成。

\textbf{正则辛形式}：
\begin{equation}
\Omega = \int_\Sigma d^3x \, \delta \mathcal{T}^\alpha_{\mu\nu}(x) \wedge \delta \pi_\alpha^{\mu\nu}(x)
\end{equation}

\textbf{基本泊松括号}：
\begin{equation}
\{\mathcal{T}^\alpha_{\mu\nu}(x), \pi_\beta^{\rho\sigma}(y)\} = \delta_\beta^\alpha \delta_{\mu\nu}^{\rho\sigma} \delta^{(3)}(x-y)
\end{equation}

其中：
\begin{equation}
\delta_{\mu\nu}^{\rho\sigma} = \frac{1}{2}(\delta_\mu^\rho \delta_\nu^\sigma - \delta_\mu^\sigma \delta_\nu^\rho)
\end{equation}

\subsubsection{Fock空间表示}

为了建立Fock空间表示，引入\textbf{相容的复结构} $J$：

\textbf{复场变量}：
\begin{equation}
\Phi_{\mu\nu}^\alpha = \frac{1}{\sqrt{2}}\left(\lambda \mathcal{T}_{\mu\nu}^\alpha + i \frac{1}{\lambda} \pi_\alpha^{\mu\nu}\right)
\end{equation}

\textbf{产生和湮灭算符}：

通过几何量子化，将经典场变量映射为算符：
\begin{align}
\hat{\mathcal{T}}_{\mu\nu}^\alpha(x) &\rightarrow \text{乘法算符} \\
\hat{\pi}_\alpha^{\mu\nu}(x) &\rightarrow -i\hbar \frac{\delta}{\delta \mathcal{T}_{\mu\nu}^\alpha(x)}
\end{align}

\textbf{正则对易关系}：
\begin{equation}
[\hat{\mathcal{T}}^\alpha_{\mu\nu}(x), \hat{\pi}_\beta^{\rho\sigma}(y)] = i\hbar \delta^\alpha_\beta \delta^{\rho\sigma}_{\mu\nu} \delta^{(3)}(x-y)
\end{equation}

%==============================================================================
\subsection{阶段B：几何量子化}
%==============================================================================
\label{subsec:stageB_geometric}

\subsubsection{预量子化线丛}

在辛流形 $(\mathcal{M}, \Omega)$ 上构造预量子化线丛 $L \to \mathcal{M}$：

\begin{itemize}
    \item 曲率：$c_1(L) = [\Omega]/2\pi\hbar$（第一陈类）
    \item 联络：$\nabla = d - i\Theta/\hbar$，其中 $\Theta$ 是辛势
    \item 截面：$\psi \in \Gamma(L)$ 是预量子化波函数
\end{itemize}

\textbf{辛势}：
\begin{equation}
\Theta = \int_\Sigma d^3x \, \pi_\alpha^{\mu\nu}(x) \, \delta \mathcal{T}_{\mu\nu}^\alpha(x)
\end{equation}

\subsubsection{Kähler极化}

选择Kähler极化，波函数为全纯函数：
\begin{equation}
\Psi[\mathcal{T}^+] = e^{-K/\hbar} f[\mathcal{T}^+], \quad K = \frac{1}{2}\int d^3x \left(\mathcal{T}^2 + \frac{\pi^2}{\mu^2}\right)
\end{equation}

其中 $K$ 是Kähler势，$\mu$ 是质量参数。

\subsubsection{与Fock空间的同构}

几何量子化的希尔伯特空间与Fock空间同构：
\begin{equation}
\mathcal{H}_{\text{geom}} \cong \mathcal{H}_{\text{Fock}} = \bigoplus_{n=0}^\infty S^n(\mathcal{H}_1)
\end{equation}

其中 $S^n$ 表示 $n$ 粒子态的对称化张量积。

%==============================================================================
\subsection{阶段C：黑洞熵的微观起源}
%==============================================================================
\label{subsec:stageC_entropy}

\subsubsection{Bekenstein-Hawking熵}

Bekenstein (1972) 和 Hawking (1974) 发现黑洞具有与视界面积成正比的熵：
\begin{equation}
S_{BH} = \frac{A}{4G} = \frac{\pi r_s^2}{G}
\end{equation}

其中 $r_s = 2GM/c^2$ 是史瓦西半径。这一公式揭示深刻联系：
\begin{itemize}
    \item 热力学：黑洞服从广义第二定律
    \item 全息：熵与面积而非体积成正比
    \item 量子：$S_{BH} \sim A/\ell_P^2$ 涉及普朗克尺度
\end{itemize}

\subsubsection{统计力学问题}

在常规热力学中，熵计算微观状态：
\begin{equation}
S = k_B \ln \Omega
\end{equation}

对于黑洞，微观自由度是什么？现有方法包括：
\begin{itemize}
    \item \textbf{弦论}：极端黑洞的D膜计数
    \item \textbf{圈量子引力}：视界处的自旋网络穿孔
    \item \textbf{诱导引力}：量子场的热熵
\end{itemize}

所有方法都需要可调参数或仅适用于特殊情况。

\subsubsection{CTUFT解决方案：谱维流动}

CTUFT从扭转场量子化提供统一的第一性原理推导。

\begin{theorem}[来自扭转场量子化的黑洞熵]
\label{thm:bh_entropy_zh}
Bekenstein-Hawking熵源于黑洞背景中扭转场的几何量子化，谱维流动 $d_s: 4 \to 10$ 提供自然的UV正则化。
\end{theorem}

\begin{proof}[推导]

\textbf{步骤1：几何量子化框架}

扭转场 $\mathcal{T}^\alpha_{\mu\nu}(x)$ 和共轭动量 $\pi_\alpha^{\mu\nu}(x)$ 构成辛流形：
\begin{equation}
\Omega = \int_\Sigma d^3x \, \delta \mathcal{T}^\alpha_{\mu\nu}(x) \wedge \delta \pi_\alpha^{\mu\nu}(x)
\end{equation}

在Kähler极化下，波函数是全纯的：
\begin{equation}
\Psi[\mathcal{T}^+] = e^{-K/\hbar} f[\mathcal{T}^+], \quad K = \frac{1}{2}\int d^3x \left(\mathcal{T}^2 + \frac{\pi^2}{\mu^2}\right)
\end{equation}

\textbf{步骤2：带谱维的态密度}

在 $d$ 维中，相对论性玻色子的态密度为：
\begin{equation}
\rho(E) \sim V \cdot E^{d/2 - 1}
\end{equation}

对于标准4D时空：$\rho_{4D}(E) \sim E^1$（体积律）。

\textbf{步骤3：视界附近的谱维流动}

在黑洞视界附近，有效谱维流动：
\begin{equation}
d_s(E) = 4 + \frac{6}{1 + (E/E_c)^2} \to 10 \quad \text{当 } E \to \infty
\end{equation}

在薄壳区域（$\rho \sim \ell_P$ 视界附近）：
\begin{equation}
\rho(E) \sim E^{10/2 - 1} = E^4
\end{equation}

\textbf{步骤4：模式计数}

能量小于 $E$ 的模式数为：
\begin{equation}
N(E) = \int_0^E dE' \, \rho(E') \sim \int_0^E dE' \, (E')^4 \sim E^5
\end{equation}

\textbf{步骤5：熵计算}

使用砖墙方法配合谱维正则化（无需人工截断）：
\begin{equation}
S = \int_0^\infty d\omega \, g(\omega) \left[\frac{\beta\omega}{e^{\beta\omega} - 1} - \ln(1 - e^{-\beta\omega})\right]
\end{equation}

对于 $d_s = 10$，$g(\omega) \sim \omega^9$，在薄壳区域 $V_{eff} \sim A \cdot \ell_P$ 积分：
\begin{equation}
S \sim \frac{A}{\ell_P^2} \cdot k_B = \frac{A}{4G}
\end{equation}

使用 $\ell_P^2 = G\hbar/c^3$ 并设 $\hbar = c = 1$。
\end{proof}

\subsubsection{关键洞见}

\begin{insight}[来自维流动的面积律]
面积律 $S \sim A$（而非 $S \sim V$）产生的原因是：
\begin{enumerate}
    \item 视界附近的高能模式经历 $d_s \approx 10$
    \item 增强的态密度 $\rho \sim E^4$ 将贡献集中在薄壳
    \item 有效体积 $V_{eff} \sim A \cdot \ell_P$ 与面积成正比
\end{enumerate}
\end{insight}

\begin{insight}[自然正则化]
与需要人工截断 $h \sim \ell_P$ 的砖墙方法不同，CTUFT通过谱维流动提供自然的UV正则化。当 $E \to \infty$ 时，$d_s \to 10$ 自动抑制发散。
\end{insight}

\subsubsection{对数修正}

完整的熵公式包含量子修正：
\begin{equation}
S = \frac{A}{4G} + \alpha \ln\left(\frac{A}{\ell_P^2}\right) + \mathcal{O}(1)
\end{equation}

\begin{theorem}[对数修正系数]
CTUFT预言 $\alpha \approx -1/2$，与圈量子引力一致。
\end{theorem}

\begin{proof}[概要]
对数修正来源于：
\begin{enumerate}
    \item $4 < d_s < 10$ 的过渡区域
    \item 视界的量子涨落
    \item 4D与10D部门之间的纠缠
\end{enumerate}

详细计算给出 $\alpha = -1/2 + \mathcal{O}(\tau_0)$，其中 $\tau_0$ 修正可忽略。
\end{proof}

\subsubsection{普适性}

\begin{theorem}[普适面积律]
熵公式 $S = A/(4G)$ 适用于所有稳态黑洞：
\begin{itemize}
    \item 史瓦西：$S = 4\pi M^2$
    \item Reissner-Nordström：$S = \pi r_+^2$，其中 $r_+ = M + \sqrt{M^2 - Q^2}$
    \item Kerr：$S = 2\pi M r_+$，其中 $r_+ = M + \sqrt{M^2 - a^2}$
    \item Kerr-Newman：$S = \pi(r_+^2 + a^2)$
\end{itemize}
\end{theorem}

\subsubsection{与其他理论的比较}

\begin{table}[h]
\centering
\caption{不同理论中的黑洞熵}
\begin{tabular}{lccc}
\toprule
\textbf{理论} & \textbf{适用范围} & \textbf{自由参数} & $\mathbf{\alpha}$ \\
\midrule
弦论 & 仅极端黑洞 & 多个 & 可变 \\
圈量子引力 & 一般 & Immirzi $\gamma$ & $-1/2$ \\
诱导引力 & 一般 & 无 & $-1$ \\
GUP模型 & 一般 & 模型依赖 & $-3/2$ 或 $-1/2$ \\
\textbf{CTUFT} & \textbf{一般} & \textbf{无}（除$\tau_0$外） & $\mathbf{-1/2}$ \\
\bottomrule
\end{tabular}
\end{table}

CTUFT优势：
\begin{itemize}
    \item 无需调谐自由参数
    \item 适用于一般（非极端）黑洞
    \item 通过谱维自然正则化
    \item 与宇宙学框架统一
\end{itemize}

%==============================================================================
\subsection{分形-扭转黑洞模型}
%==============================================================================

\subsubsection{奇点消除与向内分形展开}

\begin{hypothesis}[扭转饱和核心]
黑洞中心不是奇点，而是一个\textbf{扭转饱和的分形核心}，其中：
\begin{enumerate}
    \item 时空具有谱维 $d_s < 4$ 的分形结构
    \item 扭转达到最大临界值 $\tau_{max}$
    \item 核心具有有限密度和规则几何
\end{enumerate}
\end{hypothesis}

\begin{insight}[奇点是视角效应]
$r = 0$处的"奇点"不是无限密度的物理点。它是4D观察者对内空间展开以容纳能量的区域的视角。
\end{insight}

\begin{theorem}[规则核心定理]
在分形-扭转模型中，当 $r \to 0$ 时总应力-能量张量保持有限：
\begin{equation}
T^{total}_{\mu\nu} = T^{(matter)}_{\mu\nu} + T^{(torsion)}_{\mu\nu} < \infty
\end{equation}
其中 $T^{(torsion)}_{\mu\nu} = -\alpha(\tau^2 + \frac{1}{3}\tau^4)g_{\mu\nu}$ 提供抵消物质发散的负压。
\end{theorem}

最大密度为：
\begin{equation}
\rho_{max} = \frac{c^6}{G^3M^2}
\end{equation}

对于恒星质量黑洞（$M \sim 10M_\odot$）：$\rho_{max} \sim 10^{76}$ kg/m$^3$，虽然极高，但有限。

\subsubsection{信息悖论的解决}

\begin{insight}[信息流向内空间]
当物质落入黑洞时，它并未消失——而是过渡到10维内空间。信息被保留但对4D观察者不可达。
\end{insight}

\begin{theorem}[CTUFT中的幺正演化]
包括4D互反空间和10D内空间的总演化是幺正的：
\begin{equation}
S_{\text{total}} = S_{4D} \otimes S_{10D}, \quad S_{\text{total}}^\dagger S_{\text{total}} = \mathbb{1}
\end{equation}
\end{theorem}

\begin{proof}
完整量子态是$\mathcal{H}_{4D} \otimes \mathcal{H}_{10D}$中的矢量。黑洞形成和蒸发对应于该积空间中的幺正变换：
\begin{equation}
|\Psi_{\text{in}}\rangle_{4D} \otimes |0\rangle_{10D} \to |\text{BH}\rangle_{4D} \otimes |\Psi\rangle_{10D} \to |\Psi_{\text{out}}\rangle_{4D} \otimes |0\rangle_{10D}
\end{equation}

对4D观察者看来"丢失"的信息实际上存储在10维内空间中。
\end{proof}

\subsubsection{Page曲线}

CTUFT中霍金辐射的纠缠熵遵循Page曲线：
\begin{equation}
S_{\text{rad}}(t) = \min\left[S_{\text{BH}}(0), S_{\text{BH}}(t)\right] - \Delta S_{\text{corr}}
\end{equation}

其中$\Delta S_{\text{corr}} \sim k_B \ln d_{\text{int}}$计入内空间关联。

\subsubsection{黑洞蒸发与量子遗迹}

\begin{theorem}[黑洞终点是遗迹或爆发]
当$M \to M_{\text{crit}} \sim M_P/\tau_0$，蒸发行为改变：
\begin{itemize}
    \item 如果$\tau_0 \ll 1$：黑洞在最终爆发中完全蒸发（$\gamma$射线）
    \item 爆发能量：$E_{\text{burst}} \sim M_{\text{crit}} c^2 \sim 10^{16}$ GeV
\end{itemize}
\end{theorem}

量子遗迹：
\begin{equation}
M_{remnant} \sim M_{Planck} \cdot \tau_0^{1/3} \sim 10^{-8} \text{ kg}
\end{equation}

性质：
\begin{itemize}
    \item 质量：$M_{rem} \sim 10^{-8}$ kg
    \item 大小：$R_{rem} \sim 10^{-35}$ m（普朗克尺度）
    \item 温度：$T_{rem} \sim 10^{32}$ K
    \item 寿命：稳定（蒸发停止）
\end{itemize}

\begin{proposition}[作为暗物质的遗迹]
在早期宇宙中形成并已蒸发为遗迹的原初黑洞可能构成暗物质的显著部分。
\end{proposition}

%==============================================================================
\subsection{阶段D：原初黑洞与观测特征}
%==============================================================================
\label{subsec:stageD_primordial}

\subsubsection{原初黑洞形成}

原初黑洞（PBHs）可能形成于早期宇宙的密度涨落：
\begin{equation}
\delta \rho/\rho \gtrsim \gamma^{-1} \quad \text{在视界交叉时}
\end{equation}
其中 $\gamma \sim 0.2$ 是形成效率参数。

在时间 $t$ 形成的PBH质量为：
\begin{equation}
M_{\text{PBH}} \approx \gamma \cdot \frac{c^3 t}{G} \approx 10^{15} \left(\frac{t}{10^{-23} \text{ s}}\right) \text{ g}
\end{equation}

\subsubsection{CTUFT修正的霍金谱}

\begin{theorem}[CTUFT修正霍金辐射谱]
谱维流动修正霍金辐射谱：
\begin{equation}
n(\omega) = \frac{1}{e^{\beta\omega} - 1} \cdot \left(1 + \delta n_{d_s}(\omega)\right)
\end{equation}
其中 $\delta n_{d_s}$ 编码 $d_s \neq 4$ 效应导致的偏离：
\begin{equation}
\delta n_{d_s}(\omega) = \begin{cases}
0 & \omega \ll E_c \\
\left(\frac{\omega}{E_c}\right)^{d_s/2 - 2} - 1 & \omega \gg E_c
\end{cases}
\end{equation}
\end{theorem}

对于原初黑洞（$M \sim 10^{15}$ g），这些偏离可能在伽马射线谱的高能尾中可观测。

\subsubsection{宇宙伽马射线背景约束}

\begin{table}[h]
\centering
\caption{来自CGB的PBH丰度约束}
\begin{tabular}{lcc}
\toprule
\textbf{PBH质量 (g)} & \textbf{$f_{\text{PBH}}$约束} & \textbf{主要通道} \\
\midrule
$10^{13}$ & $\u003c 10^{-10}$ & 电子 + 光子 \\
$10^{14}$ & $\u003c 10^{-9}$ & 光子 \\
$10^{15}$ & $\u003c 10^{-8}$ & 光子 \\
$10^{16}$ & $\u003c 10^{-6}$ & 光子（次主导）
\bottomrule
\end{tabular}
\end{table}

\subsubsection{未来观测前景}

\textbf{CTA（切伦科夫望远镜阵列）}：
\begin{itemize}
    \item 能量范围：20 GeV -- 300 TeV
    \item 灵敏度：比Fermi-LAT高10--100倍
    \item 预期改进：$f_{\text{PBH}}$ 限制改善1--2个数量级
\end{itemize}

\textbf{LISA（激光干涉空间天线）}：
\begin{itemize}
    \item 对PBH双星并合敏感：$M \sim 10^{-12} - 10^{-7} M_\odot$
    \item 探测范围：对有利质量可达 $\sim$1 Gpc
    \item 预期并合率：$\sim$10--1000每年（如果 $f_{\text{PBH}} \sim 10^{-3}$）
\end{itemize}

\subsubsection{CTUFT可证伪预测}

\begin{prediction}[高能谱增强]
如果 $f_{\text{PBH}} \gtrsim 10^{-10}$ 且PBH质量 $M \sim 10^{15}$ g，CTA应在伽马射线谱的 $E \gtrsim 5 T_H \sim 50$ MeV 处观测到超出，这是 $d_s \to 10$ 效应导致的。
\end{prediction}

\begin{prediction}[PBH质量谱异常]
如果PBH构成暗物质的显著部分，LISA应探测到质量 $M \sim 10^{-12} - 10^{-7} M_\odot$ 的双星群体，其质量分布与标准天体物理形成机制不一致。
\end{prediction}

%==============================================================================
\subsection{小结：CTUFT中的黑洞}
%==============================================================================

CTUFT为黑洞物理提供完整的第一性原理描述：

\begin{enumerate}
    \item \textbf{阶段A - 正则量子化}：建立扭转场的辛结构和Fock空间表示
    \item \textbf{阶段B - 几何量子化}：构造预量子化线丛和Kähler极化
    \item \textbf{阶段C - 熵}：从扭转场量子化推导 $S = A/(4G) + \alpha\ln(A/\ell_P^2)$
    \item \textbf{阶段D - 原初黑洞}：伽马射线和引力波中的可观测特征
    \item \textbf{奇点}：被有限密度扭转饱和核心取代
    \item \textbf{信息}：在10维内空间中保留
\end{enumerate}

关键预测：
\begin{itemize}
    \item 对数修正 $\alpha \approx -1/2$
    \item 原初黑洞的修正霍金谱（$d_s \to 10$ 增强）
    \item PBH丰度限制：$f_{\text{PBH}} \lesssim 10^{-8}$（$M \sim 10^{15}$ g）
    \item 量子遗迹作为暗物质候选者
    \item 如果 $f_{\text{PBH}} \gtrsim 10^{-10}$，CTA可探测伽马射线超出
    \item 如果 $f_{\text{PBH}} \gtrsim 10^{-3}$，LISA可探测PBH双星
\end{itemize}

理论推导（阶段A、B、C）与观测约束（阶段D）的相互作用为通过黑洞物理检验CTUFT提供了全面框架。
